\subsection{Vector Memory Operations}

A vector memory operand uses the (:term:base.terms.ordinary:) compact, \texttt{EXT1}, and
\texttt{EXT2} descriptor formats in vector context.  The descriptor computes
one scalar anchor address.  A contiguous access uses

\[
  address[i] = anchor + iE,
\]

All arithmetic is modulo $2^{64}$. Existing effective-address auto-update is
applied once to the scalar anchor in the (:term:base.terms.ordinary:) architectural order. That
single update covers all lanes.

Inactive lanes issue no memory transaction.  (:ref:VECTOR.instructions.VMOV:) merges a memory
source into its destination, while (:ref:VECTOR.instructions.VMOVZ:) writes zero to inactive
destination lanes. A memory-destination arithmetic form uses non-atomic,
lane-wise read-modify-write semantics. A vector instruction nevertheless
has one architectural commit point: address evaluation, all active reads,
computation, and all active writes must succeed before any register,
predicate, memory, auto-update, \texttt{FFLAGS}, or PC effect becomes visible.
After range, wrap, and canonicality checks, the implementation completes
translation of every mapping covered by every active lane in increasing
linear-address order.  Only after all such translations succeed does it apply
physical-class checks, the vector/gather/scatter MMIO-operation prohibition,
alignment checks, and target or bus accesses, in that order.  Consequently, a
translation fault in a higher-address active lane takes precedence over an
MMIO condition discovered in a lower-address active lane.  When logical lane
order differs from address order, translation and translation-fault selection
still follow address order.  The reported fault address identifies the lane
whose range supplies the selected failure.
